\section{Experimental Design}\label{sec:setup}

\subsection{Test bench}
The bench has four parts that share one TypeScript package, \texttt{sync-core}, holding the data model, the
operations, the clocks, the CRDT building blocks and the five strategies behind the interface of
Definition~\ref{def:strategy}.
\begin{itemize}
\item The \emph{application}, SubTrack, is a React Native subscription tracker that renders $\mathit{read}$, lets
  the user resolve $\mathit{conflicts}$, persists the replica in SQLite and has a developer screen for the strategy,
  clock skew, forced offline mode and metric export.
\item The \emph{server} is a Node.js process (a Docker image) that keeps its replica in SQLite, answers one
  synchronisation round per HTTP request and notifies clients through a WebSocket.
\item The \emph{simulator} is a deterministic discrete-event program that runs $n$ clients and the server in one
  process, with virtual time, seeded pseudo-random streams, online and offline periods and per-client clock skew. It
  generates the user intents from the causal history of each client, so every strategy receives exactly the same
  operations, and it records $\mathit{seen}$ for every operation.
\item The \emph{oracle} applies Definitions~\ref{def:outcome} and~\ref{def:metrics} to the final state.
\end{itemize}
Simulated runs execute the same strategy code as the application.

\subsection{Choice of technologies}
Table~\ref{tab:tech} lists the choices and the alternatives considered. The criteria were that one implementation of
each strategy must run unchanged on the phone, on the server and in the simulator, and that the CRDT baselines must be
libraries used in practice.

\begin{table}[ht]
\caption{Technology choices and the alternatives that were not chosen.}
\label{tab:tech}
\footnotesize
\begin{tabular}{@{}>{\raggedright\arraybackslash}p{0.2\textwidth}>{\raggedright\arraybackslash}p{0.75\textwidth}@{}}
\toprule
Chosen & Alternatives and reason \\
\midrule
React Native, Expo & Flutter or native Kotlin/Swift would need a second implementation of every strategy; with
  TypeScript one implementation runs on the device, the server and in the simulator. \\
SQLite on the device & Realm and WatermelonDB bring their own sync model, PouchDB/CouchDB its revision-tree conflict
  handling; a neutral store keeps the strategy the only variable. \\
Node.js, SQLite server & Same language as the strategies; a single-file database is enough for one hub and is easy
  to reset between runs. \\
Automerge, Yjs & The two most used CRDT libraries, with different designs; a hand-written full CRDT would measure
  our code, not the state of practice. \\
In-process simulation & Device farms give realistic timing but cannot reproduce a run or control clocks and
  connectivity; devices only validate a subset (Section~\ref{sec:validation}). \\
\bottomrule
\end{tabular}
\end{table}

One constraint was found during the integration: the Automerge core is compiled to WebAssembly, which Hermes, the
JavaScript engine of React Native, does not execute. The application therefore ships S1, S2, Yjs and S4, and
Automerge is evaluated in the simulator and on the server, where Node.js runs it.

\subsection{Experiments}
Table~\ref{tab:experiments} lists the experiments. Each varies one or two independent variables of the simulator
while the others stay at the defaults of the random workload of Section~\ref{sec:case}.

\begin{table}[ht]
\caption{Planned experiments (S3 stands for both Automerge and Yjs).}
\label{tab:experiments}
\footnotesize
\begin{tabular}{@{}l>{\raggedright\arraybackslash}p{0.6\textwidth}ll@{}}
\toprule
Exp. & Measure and independent variables & Strategies & RQ \\
\midrule
E1 & loss vs.\ number of clients $n \in \{2,5,10,20\}$, offline duration 1~min--24~h, edit rate & S1--S3 & RQ1 \\
E2 & loss vs.\ clock skew $0$, $\pm 1$~s, $\pm 1$~min, $\pm 1$~h, $\pm 1$~day & S1, S2, S4 & RQ1 \\
E3 & payload, state, sync time vs.\ $10^2$--$10^4$ operations, $10$--$10^3$ records & S1--S3 & RQ2 \\
E4 & loss and cost per field class: SIMPLE, SET, COUNTER only, mixed & S1--S3 & RQ3 \\
E5 & the hybrid on the workloads of E1--E4 & S1--S4 & RQ4 \\
E6 & implementation complexity ($\mathit{loc}$) & S1--S4 & RQ2, 4 \\
\bottomrule
\end{tabular}
\end{table}

\subsection{Validation method}
Each configuration is run with 30 seeds. We report the mean with a 95\% confidence interval based on Student's
$t$ distribution. Since loss rates and byte counts are not normally distributed, pairs of strategies are compared with
the two-sided Mann--Whitney $U$ test~\cite{mann1947test}, and the $p$-values of all comparisons of one experiment are
adjusted with the Holm procedure~\cite{holm1979simple}. S4 meets RQ4 if, on E1--E4, it loses significantly fewer
operations than S1 ($p < 0.05$ after correction), converges in every run, and has a significantly smaller mean payload
and client state than S3; a partial result is reported as such, with the trade-off observed.

\subsection{Validation on public data and on devices}\label{sec:validation}
Two checks guard against artefacts of the bench. First, the measurement harness is validated on public data: the
workloads of the crdt-benchmarks suite~\cite{crdtbenchmarks}, including its trace of a real editing session, are
replayed through our Automerge and Yjs adapters, and the encoded sizes are compared with the values published for the
same library versions; a large difference would point to a fault in our adapters or byte counting. Second, a subset of
the scenarios of Section~\ref{sec:case} is replayed by hand on two Android emulators and one phone against the Docker
server, using the developer screen, and the outcomes and byte counts are compared with those of the simulator.
